% !TEX program = pdflatex
\documentclass[11pt,a4paper]{article}

\usepackage[margin=1in]{geometry}
\usepackage{amsmath,amssymb,amsthm,mathtools}
\usepackage[T1]{fontenc}
\usepackage{lmodern}
\usepackage{microtype}
\usepackage{booktabs}
\usepackage{array}
\usepackage{xcolor}
\usepackage{hyperref}
\usepackage{enumitem}
\usepackage{listings}
\usepackage{longtable}

\hypersetup{
    colorlinks=true,
    linkcolor=blue,
    citecolor=blue,
    urlcolor=blue,
    pdftitle={2-Adic Branch Collapse in Odd Semiprime Factorization Experiments},
    pdfauthor={Research Log Synthesis}
}

\newtheorem{definition}{Definition}[section]
\newtheorem{proposition}[definition]{Proposition}
\newtheorem{conjecture}[definition]{Conjecture}
\newtheorem{remark}[definition]{Remark}

\newcommand{\vTwo}{v_2}
\newcommand{\N}{\mathbb{N}}
\newcommand{\Z}{\mathbb{Z}}

\title{\textbf{2-Adic Branch Collapse in Odd Semiprime Factorization Experiments}\\
\large A Synthesis of the Structural Results Leading to Experiment 648}
\author{Research Log Synthesis}
\date{21 August 2026}

\begin{document}
\maketitle

\begin{abstract}
This paper synthesizes a sequence of computational experiments investigating the arithmetic structure exposed by repeatedly reducing an odd semiprime modulo powers of two. The central object is an odd semiprime
\[
    n=pq,
\]
with odd prime factors $p\le q$. The experiments began by recording residue information of $n$ and the associated factor branches modulo $2^k$, then progressively compressed that information into valuations, divisor-level collisions, and frame-dependent branch coordinates. The resulting picture is substantially simpler than the original residue tree suggests.

For the two observed affine frames, the relevant branch coordinates are
\[
    \text{Frame A:}\qquad A=p-3,\quad B=q+3,\quad n+9=AB,
\]
and
\[
    \text{Frame B:}\qquad A=p+1,\quad B=q-3,\quad n+3=AB.
\]
Because $A$ and $B$ arise from odd-prime factors and are generally even, the depth of the common dyadic branch is governed by the minimum of their $2$-adic valuations. The experiments indicate the exact residual rule
\[
    m=\min\{\vTwo(A),\vTwo(B)\},
\]
and then test whether the second branch can itself be eliminated using
\[
    \vTwo(n+c),
\]
leading to the candidate collapse
\[
    m=\min\{\vTwo(p-p_0),\\vTwo(n+c)\},
\qquad
    d=\min\{\vTwo(p-p_0)+1,\vTwo(n+c)\}.
\]
Here $(p_0,c)=(3,9)$ in Frame A and $(p_0,c)=(-1,3)$ in Frame B.

The paper distinguishes identities derived directly from the algebra from statements supported only by finite computation. No claim of a general polynomial-time factoring algorithm is made. The main achievement is a structural compression: a large residue signature can, within the tested family, collapse to a small valuation signature that describes the branch depth exactly.
\end{abstract}

\section{Introduction}

A recurring theme of the research was that an odd semiprime can contain much more arithmetic structure than is visible from its numerical magnitude alone. Rather than attempting to factor
\[
    n=pq
\]
directly, the experiments examined what can be learned by observing $n$ modulo
\[
    2,4,8,16,\ldots,2^k.
\]

At first this naturally produces a binary lifting tree. Each increase in $k$ reveals another bit of the unknown factors. Early experiments treated these residue patterns explicitly, and later experiments searched for lower-dimensional descriptions of the same information. The progression was approximately:

\begin{enumerate}[leftmargin=2em]
    \item residue lifting modulo powers of two;
    \item reconstruction and synchronization of factor branches;
    \item affine changes of variables that expose the repeated constants;
    \item divisor-level collision structure of shifted products $n+c$;
    \item reduction of branch depth to $2$-adic valuations;
    \item testing whether the two branch valuations collapse to one factor-side valuation plus one $n$-only valuation.
\end{enumerate}

The current paper records that progression and states the mathematical structure that has emerged by Experiment 648.

\section{Problem Setting}

Throughout, let
\[
    n=pq,
\]
where $p$ and $q$ are odd primes. The experiments do not assume that the prime factorization is known to the algorithm under investigation; the known factorization is used only by the experimental harness to verify the proposed laws.

For an integer $x\ne0$, define the $2$-adic valuation
\[
    \vTwo(x)=\max\{r\ge0:2^r\mid x\}.
\]
We use the convention $\vTwo(0)=+\infty$ when needed.

The computational question is not merely whether residues distinguish factors. The stronger question is:

\begin{quote}
What is the smallest invariant that determines the observed branch depth?
\end{quote}

This distinction is crucial. A complete residue signature may determine an answer without being the essential reason the answer has its particular form.

\section{From Residue Trees to Branch Coordinates}

For an odd semiprime, the residue class of $n$ modulo $4$ is especially informative:
\[
    n\equiv1\pmod4
    \quad\text{or}\quad
    n\equiv3\pmod4.
\]
In the experiments this split was used to select two affine coordinate systems. They are most naturally described as follows.

\subsection{Frame A}

For the branch class corresponding to $n\equiv3\pmod4$, define
\[
    A=p-3,
    \qquad
    B=q+3.
\]
Then
\[
\begin{aligned}
    AB
    &=(p-3)(q+3)\\
    &=pq+3p-3q-9.
\end{aligned}
\]
The structural experiments associate this frame with the shifted product
\[
    n+9,
\]
so that the branch calculation is organized around
\[
    \vTwo(p-3),\qquad \vTwo(q+3),\qquad \vTwo(n+9).
\]

\subsection{Frame B}

For the complementary branch class, define
\[
    A=p+1,
    \qquad
    B=q-3.
\]
The corresponding shifted quantity used in the experiments is
\[
    n+3.
\]
Thus the key valuation variables become
\[
    \vTwo(p+1),\qquad \vTwo(q-3),\qquad \vTwo(n+3).
\]

The important observation is not the particular constants by themselves, but the existence of a frame in which the branch offsets become fixed affine displacements from small odd anchors.

\section{The Basic 2-Adic Mechanism}

The relevant arithmetic phenomenon is elementary but powerful. Suppose two even integers $A$ and $B$ have distinct $2$-adic valuations:
\[
    \vTwo(A)<\vTwo(B).
\]
Then
\[
    \vTwo(A+B)=\vTwo(A).
\]
The same statement holds with $A-B$ in place of $A+B$.

Therefore, for a product relation of the form
\[
    n+c=AB,
\]
the valuation satisfies
\[
    \vTwo(n+c)
    =
    \vTwo(A)+\vTwo(B),
\]
provided the displayed equality is an exact integer identity. More generally, when the branch construction is based on a sum/difference relation rather than a literal product identity, the smaller residual valuation controls the first level at which the branches cease to be indistinguishable. This is the source of the observed ``minimum valuation'' law.

The computational experiments repeatedly encountered the same dichotomy:

\begin{enumerate}[leftmargin=2em]
    \item if the two residual valuations are unequal, the smaller one controls the collision depth;
    \item if the two residual valuations are equal, one additional $2$-adic level survives before the branches separate.
\end{enumerate}

That second case is the origin of the familiar extra-one correction in the depth formula.

\section{The Residual-Valuation Law}

Let
\[
    a=\vTwo(A),
    \qquad
    b=\vTwo(B).
\]
The experiments identify the collision depth parameter $m$ with
\[
    \boxed{m=\min(a,b)}.
\]

The corresponding exact-depth rule is
\[
    \boxed{
    d=m+\mathbf{1}_{a=b}
    },
\]
where $\mathbf{1}_{a=b}$ is $1$ when $a=b$ and $0$ otherwise.

This can be rewritten as
\[
    \boxed{
    d=\min(a+1,b+1)
    }
\]
for the equal-valuation branch and, more usefully for the observed experiments,
\[
    d=\min(a+1,w)
\]
once the $n$-only valuation $w=\vTwo(n+c)$ is known to encode the same collision boundary.

The key conceptual step is that a branch tree that appears to require the complete sequence of residues modulo $2^k$ can instead be described by a small number of integers.

\section{Experiment 614 and the Shifted-Divisor View}

A major transition occurred when the analysis moved from raw residue classes to divisors of shifted quantities $n+c$. Instead of asking independently what happens modulo every power of two, the experiments examined the divisor structure generated by a small affine shift of $n$.

This changed the question from
\[
    \text{Which residue branch survives at level }k?
\]
to
\[
    \text{Which dyadic divisor is shared by the relevant shifted factors?}
\]

The latter is naturally measured by a $2$-adic valuation. A dyadic common divisor
\[
    2^m\mid A,\qquad 2^m\mid B
\]
has largest possible exponent
\[
    m=\min\{\vTwo(A),\vTwo(B)\}.
\]

The experiments following this transition therefore stopped treating every bit as an independent empirical feature. The bits became manifestations of a single divisor exponent.

\section{Experiment 635: Divisor-Level Collision Structure}

Experiment 635 formalized this change of viewpoint. The computational record was explicitly organized around the divisor-level collision structure of $n+c$ over a large semiprime sample.

The important lesson was that the collision was not an arbitrary statistical pattern. It aligned with divisibility. When two candidate branches remained compatible modulo $2^k$, that compatibility could be represented as membership in a dyadic divisor class.

Thus the effective invariant was already moving toward
\[
    m=\vTwo(\text{relevant branch displacement}).
\]

The experiment did not by itself prove the final one-valuation formula, but it supplied the divisor-theoretic bridge from a residue tree to a valuation problem.

\section{Experiments 636--647: Frame Structure and Exact Branches}

The subsequent experiments refined the branch coordinates. The principal discovery was that the two residue branches could be described using fixed affine offsets tied to the parity frame.

The resulting coordinate systems are:

\begin{center}
\begin{tabular}{lll}
\toprule
Frame & factor-side displacement & shifted quantity \\
\midrule
A & $p-3$, $q+3$ & $n+9$ \\
B & $p+1$, $q-3$ & $n+3$ \\
\bottomrule
\end{tabular}
\end{center}

Experiment 647 then established the strongest version reached before the present work: for the tested domain, the full residue signature
\[
    (\text{frame},\;n\bmod 2^k,\;p\bmod 2^k)
\]
determined the observed collision quantity $m_k$ exactly.

That result was important, but it still left a natural question. The residue of $p$ modulo $2^k$ contains many bits. Was all of that information truly necessary?

Experiment 648 was designed specifically to answer that question.

\section{Experiment 648: Single-Valuation Branch Collapse}

Experiment 648 tests whether the complete branch information can be compressed to one factor-side valuation together with one valuation computed entirely from $n$.

For Frame A define
\[
    p_0=3,
    \qquad c=9,
\]
and for Frame B define
\[
    p_0=-1,
    \qquad c=3.
\]
Then put
\[
    t=\vTwo(p-p_0),
    \qquad
    w=\vTwo(n+c).
\]

The proposed collapse is
\[
    \boxed{m=\min(t,w)}
\]
and therefore
\[
    \boxed{d=\min(t+1,w)}.
\]

Written by frame, the candidate formula becomes
\[
    \boxed{
    d_A=
    \min\left(\vTwo(p-3)+1,\;\vTwo(n+9)\right)
    }
\]
and
\[
    \boxed{
    d_B=
    \min\left(\vTwo(p+1)+1,\;\vTwo(n+3)\right)
    }.
\]

This is the central compression hypothesis of the current stage of the research.

\section{Why the Collapse Is Plausible}

Suppose the relevant two residual valuations are
\[
    a=\vTwo(A),
    \qquad
    b=\vTwo(B).
\]
Let
\[
    m=\min(a,b).
\]
If $a\ne b$, then the lower valuation dominates the first nonzero dyadic contribution, so the shifted quantity records the same boundary:
\[
    w=m.
\]
If $a=b$, cancellation survives one additional level, producing
\[
    w>m.
\]
Hence the shifted valuation distinguishes the two regimes:

\[
\begin{array}{c|c|c}
\text{case} & w & d \\
\hline
 a\ne b & w=m & m \\
 a=b & w>m & m+1
\end{array}
\]

At the level of the proposed one-sided description, this is exactly
\[
    d=\min(a+1,w).
\]

Thus $w$ plays a dual role: it measures the available dyadic depth of the shifted integer and simultaneously detects whether equal residual valuations allow the extra branch level.

\section{A General Valuation Lemma Behind the Pattern}

The observed behavior is consistent with the following elementary principle.

\begin{proposition}[Unequal valuation dominance]
Let $x,y\in\Z\setminus\{0\}$. If
\[
    \vTwo(x)<\vTwo(y),
\]
then
\[
    \vTwo(x+y)=\vTwo(x)
\]
and
\[
    \vTwo(x-y)=\vTwo(x).
\]
\end{proposition}

\begin{proof}
Write
\[
    x=2^a u,
    \qquad
    y=2^b v,
\]
where $u,v$ are odd and $a<b$. Then
\[
    x\pm y
    =2^a\left(u\pm2^{b-a}v\right).
\]
Since $b-a\ge1$, the second term in parentheses is even while $u$ is odd. Hence the parenthesis is odd, giving valuation $a$.
\end{proof}

When $a=b$, the odd parts can cancel modulo an additional power of two, which is why equality of valuations is the exceptional case in the branch-depth law.

\section{What Has Been Learned}

The sequence of experiments yields several layers of compression.

\subsection{The residue tree is structured, not arbitrary}

The successive values of $n\bmod 2^k$ do not produce unrelated observations. They are nested constraints generated by the same $2$-adic structure.

\subsection{The factor branches have affine anchors}

The repeated offsets $3$, $-3$, and $1$ are not merely empirical constants in the current representation. They define useful affine coordinates in which the parity branches become simple valuation questions.

\subsection{Divisibility is more fundamental than the raw bit string}

The observable branch depth is naturally described by the largest power of two dividing the relevant residual displacement. Thus a long residue history can encode a single integer valuation.

\subsection{The two-factor state may collapse to one side plus $n$}

The most important current reduction is the candidate replacement
\[
    (\vTwo(A),\vTwo(B))
    \quad\longrightarrow\quad
    \left(\vTwo(p-p_0),\vTwo(n+c)\right).
\]
If exact beyond the tested range, this would eliminate one factor-side valuation from the state description.

\subsection{Depth and recoverability are different questions}

The experiments establish structural relations among residues, valuations, and branch depths. They do not automatically imply that the unknown prime factor $p$ can be reconstructed efficiently from $n$ alone. A formula containing
\[
    \vTwo(p-p_0)
\]
still contains information about an unknown factor.

This distinction prevents an overstatement of the result. The experiments uncover a compressed description of the factor-dependent branch state; they do not yet provide an unconditional factoring algorithm.

\section{Computational Methodology}

The experiments use exact integer arithmetic. A representative implementation performs the following steps:

\begin{enumerate}[leftmargin=2em]
    \item generate odd primes up to a fixed limit;
    \item form semiprimes $n=pq$ with $p\le q$;
    \item select Frame A or Frame B from $n\bmod4$;
    \item compute the exact branch coordinates;
    \item compute $2$-adic valuations using integer arithmetic;
    \item compare observed collision depths with the candidate formulas;
    \item group states by compressed signatures and search for collisions.
\end{enumerate}

The central validation criterion is exact equality, not correlation. A candidate law fails if even a single tested semiprime produces a different predicted value.

In particular, Experiment 648 checks:
\[
    m
    \stackrel{?}{=}
    \min\bigl(\vTwo(p-p_0),\vTwo(n+c)\bigr),
\]
\[
    m
    \stackrel{?}{=}
    \min\bigl(\vTwo(q-q_0),\vTwo(n+c)\bigr),
\]
and
\[
    d
    \stackrel{?}{=}
    \min\bigl(\vTwo(p-p_0)+1,\vTwo(n+c)\bigr).
\]

A particularly strong diagnostic is signature injectivity. If two test instances share the same compressed signature
\[
    (\text{frame},\vTwo(p-p_0),\vTwo(n+c)),
\]
then they must also have the same observed $m$ and $d$. Any counterexample would show that the proposed compression discarded information that the depth actually depends on.

\section{Interpretation of the Frames}

The two frames can be viewed as local $2$-adic coordinate charts.

\subsection{Frame A}

The anchor $p=3$ is singled out, and the relevant displacement is
\[
    p-3.
\]
Large values of
\[
    \vTwo(p-3)
\]
mean that $p$ lies unusually close to $3$ in the $2$-adic metric. The farther $p$ is from $3$ in the $2$-adic sense, the deeper the branch can remain synchronized before that branch displacement becomes visible.

\subsection{Frame B}

The corresponding anchor is $p=-1$ modulo powers of two, represented by the displacement
\[
    p+1.
\]
Again,
\[
    \vTwo(p+1)
\]
measures $2$-adic proximity to the frame anchor.

This gives a useful geometric interpretation: the branch depth is controlled by how close the factor lies to a small affine anchor in the $2$-adic metric, subject to the ceiling imposed by the shifted value $n+c$.

\section{A Unified Notation}

For compactness, define the frame-dependent pair
\[
    (p_0,c)=
    \begin{cases}
        (3,9), & \text{Frame A},\\
        (-1,3), & \text{Frame B}.
    \end{cases}
\]
Then the candidate law is simply
\[
    \boxed{
    d=
    \min\left(
        \vTwo(p-p_0)+1,
        \vTwo(n+c)
    \right).
    }
\]

Equivalently, with
\[
    t=\vTwo(p-p_0),
    \qquad
    w=\vTwo(n+c),
\]
we have
\[
    d=
    \begin{cases}
        w, & w\le t,\\
        t+1, & w>t.
    \end{cases}
\]

This form makes the limiting mechanism explicit: either the shifted integer $n+c$ supplies the limiting dyadic depth, or the factor's distance from the frame anchor supplies it and allows one additional level.

\section{Limits and Unresolved Questions}

Several questions remain open.

First, the present results are computationally verified over finite ranges. Exact finite verification can reject a conjecture on the tested domain, but cannot by itself establish universal validity.

Second, the frame selection currently depends on the parity structure of the semiprime. A complete theory should derive the frames from first principles rather than treating them as discovered empirical coordinates.

Third, the quantity
\[
    \vTwo(p-p_0)
\]
is still factor-dependent. To obtain a genuinely factor-free predictor from $n$ alone, one would need either a way to recover this valuation directly from observable data or a theorem proving that it is itself determined by a smaller invariant of $n$.

Fourth, it remains to determine how the discovered $2$-adic structure interacts with odd prime constraints beyond the tested ranges. The fact that $p$ and $q$ are prime may impose additional restrictions that have not yet been fully exploited.

Finally, the relationship between this branch-depth structure and actual factor recovery remains unresolved. A structural invariant can be exact while still failing to identify the factors uniquely.

\section{Research Direction Suggested by Experiment 648}

If the single-valuation collapse survives substantially larger tests, the next experiments should focus on eliminating variables rather than adding more features.

The natural progression is:

\begin{enumerate}[leftmargin=2em]
    \item verify the formula over several independent prime ranges;
    \item test all admissible orderings $p\le q$ and both frames;
    \item search for an $n$-only expression for $\vTwo(p-p_0)$;
    \item classify all collisions under the compressed signature;
    \item determine whether the branch depth identifies a unique factor residue class;
    \item investigate whether the valuation law extends from $2$-adic depth to an explicit reconstruction algorithm.
\end{enumerate}

A particularly important test is to determine whether two distinct factor pairs with the same $n$, modulo the observable information, can share the same compressed depth. If so, the depth is a structural invariant but not a complete factor identifier. That would still be mathematically valuable, because it cleanly separates ``branch geometry'' from ``factor recovery.''

\section{Conclusion}

The central lesson of the experiments is that the apparent complexity of the repeated-modulus factor branches is misleading. What begins as a growing signature
\[
    n\bmod2,
    n\bmod4,
    n\bmod8,
    \ldots
\]
can be reorganized into a small number of affine $2$-adic quantities.

The strongest structural picture reached so far is summarized by the frame-dependent coordinates
\[
    \text{A: }(p-3,\;q+3,\;n+9),
\]
\[
    \text{B: }(p+1,\;q-3,\;n+3),
\]
and by the valuation law
\[
    m=\min\{\vTwo(A),\vTwo(B)\}.
\]

Experiment 648 asks whether this can be compressed one step further. The candidate is
\[
    m=\min\left(\vTwo(p-p_0),\vTwo(n+c)\right),
\]
which yields
\[
    \boxed{
    d=\min\left(\vTwo(p-p_0)+1,\vTwo(n+c)\right).
    }
\]

If this identity is confirmed generally, the research has moved from a residue-tree description to a valuation-level description and, potentially, from a two-branch factor state to a one-sided factor valuation plus an $n$-only ceiling.

That is the main conceptual result of the current stage: the binary lifting process appears to be governed not by an ever-growing collection of independent residue bits, but by a small $2$-adic geometry centered on a handful of affine anchors.

\section*{Appendix: Minimal Experiment 648 Core}

The essential computational test can be reduced to the following pseudocode.

\begin{lstlisting}[basicstyle=\ttfamily\small,breaklines=true]
for each odd semiprime n = p*q:
    if n mod 4 == 3:
        p0 = 3
        c  = 9
        A  = p - 3
        B  = q + 3
    else:
        p0 = -1
        c  = 3
        A  = p + 1
        B  = q - 3

    a = v2(A)
    b = v2(B)
    m = min(a, b)
    d = m + (a == b)

    t = v2(p - p0)
    w = v2(n + c)

    assert m == min(t, w)
    assert d == min(t + 1, w)
\end{lstlisting}

The assertions are exact integer identities on the tested sample. Their mathematical status remains empirical until a proof is supplied.

\end{document}
